\section{Possibility premises and fixed-state reasoning}
\label{sec:annex-p19-conditional-interpretation-experiments}

\begin{annexitem}{annex:p19.fixed-state-mp}{thm:p19-fixed-state-mp}{Restriction, order and the compatibility premise}
\begin{proof}
If \(I\) accepts \(p\), every \(v\in I\) lies in \([p]\).
The condition \(I\cap[p]\subseteq[q]\) therefore puts every such
\(v\) in \([q]\). This is acceptance of \(q\). A further existential
condition about \(r\) removes none of these hypotheses and cannot change
this implication.

Conditional exclusion followed by accepting \(p\) gives
\[
 \bigl(I\setminus([p]\cap[\neg q])\bigr)\cap[p]
 =I\cap[p]\cap[q].
\]
Indeed, membership on the left requires membership in \(I\), truth of
\(p\), and failure of \(p\land\neg q\); with \(p\) true this last
condition is exactly truth of \(q\). In the opposite order the result is
\[
 (I\cap[p])\setminus([p]\cap[\neg q])=I\cap[p]\cap[q]
\]
by the same membership argument. If this set is nonempty it is a state
all of whose valuations make \(q\) true. If it is empty, the two accepted
constraints have no valuation in the original state, and the statement
does not relabel the empty set a nonempty state.

Suppose now that both restricted conditionals hold, and assume for a
contradiction that \(v\in I\cap[p]\cap[r]\). The first conditional makes
\(q\) true at \(v\); the second makes \(\neg q\) true there. The Boolean
negation clause forbids this. Consequently no such \(v\) exists, proving
incompatibility with \(\operatorname{Poss}_I(p\land r)\). The state
\(I=\{(1,1,0)\}\) shows why the compatibility hypothesis matters: the
first conditional holds with its antecedent true, and the second holds
because no valuation in \(I\) satisfies its antecedent. Thus there is no
unqualified inconsistency claim about the two conditionals alone.
For the final claim take this same \(I=\{(1,1,0)\}\) and the valuation
\(w=(1,0,1)\) outside it. The restricted conditional is true on \(I\),
whereas \(w\) makes \(p\) true and \(q\) false, so the material
conditional is false at \(w\). No premise put \(w\) in \(I\).
\end{proof}
\end{annexitem}

\begin{annexitem}{annex:p19.might-expansion}{thm:p19-might-expansion}{Every valuation in the two expansions}
\begin{proof}
In \(I_0\), its first valuation has \(p=0\), and its second has
\(p=q=1\). Hence there is no \(p\land\neg q\) valuation and
\(C_{I_0}(p,q)\) holds. Both have \(r=0\), so \(r\) is not possible.
Each proposed expansion adds a valuation with \(r=1\), and therefore
does make \(r\) possible. Restricting to \([p]\) removes precisely
\((0,0,0)\) and yields the two final states displayed in the body.

Their remaining valuations evaluate the relevant formulas as follows:
\[
\begin{array}{c|ccccc}
 (p,q,r)&p&q&\neg q&r&r\lor q\\\hline
 (1,1,0)&1&1&0&0&1\\
 (1,1,1)&1&1&0&1&1\\
 (1,0,1)&1&0&1&1&1
\end{array}
\]
The wet state uses rows one and two; the dry state uses rows one and
three. In both, all rows satisfy \(p\) and \(r\lor q\), some row
satisfies \(r\), and row one falsifies \(r\). Thus compatibility with
\(r\) is not acceptance of \(r\). Both wet rows satisfy \(q\); in the
dry state row three falsifies \(q\) and row one falsifies \(\neg q\).

Because every final row has \(p=1\), the restricted conditional
\(C_I(p,q)\) holds exactly when every final row has \(q=1\). It holds
on the wet state and fails on the dry one. This proves that accommodating
the same existential condition can preserve or destroy this conditional,
depending on the chosen expansion. No unique expansion was included in
the definition of compatibility, so neither choice follows from that
definition alone. Finally, intersecting the dry state with \([\neg r]\)
retains row one and removes row three. The resulting nonempty state
accepts \(q\), as claimed. Every conclusion follows from the stipulated
valuations; none uses a psychological frequency as a premise.
\end{proof}
\end{annexitem}
